\section{A simple account: two cards, one slot}
\label{sec:account}

A simple mechanism could produce most of these results. Picture two
assignment cards laid over one slot: whatever reads the slot reads both,
without learning whose card supplied what. In the model, the slot is the
stretch of positions where A's own card lies, B's card is laid over it,
and the reader is each attention head, which returns one blend of both.
Three properties of the memory link give this picture a concrete
basis. First, by Eq.~\ref{eq:bridge}, each
head returns $o=(1-\beta)\,o_{\mathrm{own}}+\beta\,o_B$, a blend of A's
own value average and B's, weighted by $\beta$, the share of attention
that falls on B's keys. The blend leaves the head as a single vector,
with no field saying which part came from which cache. A could still
tell the sources apart if B's card stated its assignment in a different
place or layout from A's own, and the next two properties show that it
does not. Second, B's keys keep their original rotary positions, so,
seen from A's queries, B's card occupies nearly the same positions as
A's own. Third, both cards follow one template. Where A's own card says
\qq{You were assigned the priority \ldots}, B's card states its
assignment in the same place, within the same stretch of positions.
B's card does name B, but the account does not need A to read names.
It adds one premise, that A finds its own assignment by where it lies, reading
whatever occupies that stretch of its context as its own. The three
properties guarantee that B's assignment occupies it too.

The account fits results we already had, so the fit is weak evidence;
what could refute it are the predictions at the end of this section. Its
strongest support comes from two results that the obvious alternatives
do not predict.

If A claimed B's rule because B's card addresses its reader as
\qq{you}, a third-person record that names B wherever it states B's
assignment should have stopped the claim. On the account the name should
not matter, because it changes what B's card says but not where the card
lies. It did not matter: A claimed B's rule from the third-person record
about as often as from the original card (Section~\ref{sec:third}).

If claiming instead followed how much of B's content reached A, then
under a steering vector tuned to give A the same access to B's rule
as the link at $w=1$, A should have claimed B's rule as often as
under the link. It never did (Section~\ref{sec:route}). On the account
this is expected, because the vector is a direction added to A's hidden
states, not a card that A's queries can read as an assignment.

The remaining results, including the one text header that reversed our
prediction, fit the account as well, though they discriminate less
(Appendix~\ref{app:furtherdisc}). One result lies outside it: although
the blend in each head is graded, single answers were sharp
(Section~\ref{sec:sharp}), and the account does not say how a blend
becomes a choice.

The account makes testable predictions; the sharpest could refute it.
Place B's cache at positions that do not overlap A's card, as
prepended-cache systems such as LatentMAS do \citep{zou2026latentmas},
and give it no header. The account predicts claiming near the untagged
message's level, far below our memory link's; if claiming stays at the
memory link's level, the account is wrong. Three further predictions
follow. Claiming should fall if B's card is laid out unlike A's, for
example with its assignment in another place or in a table, because
B's assignment would then no longer lie where A's own lies. The
third-person record did not do this: it changed the grammatical person
of the assignment line but not where the line lies. Giving B's card a different
entry for Robin should make A's answers about Robin move to B's version,
as its answers about itself moved to B's rule. Finally, deleting B's
code word from A's connected reflection should stop that word from
persisting after the cut, which would turn the post hoc grouping of
Section~\ref{sec:moments} into a causal test.

\section{Discussion}
\label{sec:discussion}

\paragraph{The self as a default.} Two regularities run through our
results. Content laid over the receiver's own card became the receiver's
own, even when the content itself named its owner, as B's card and the
third-person record both did. The line between self and other held most
firmly where the source was named apart from the content, in the header
of a message that did not lie over the receiver's card. Moving the
content off the receiver's card without naming a source already cut
claiming from 95.8\% to 14.6\%, though position, form, and the content
itself (B's reflection rather than B's card) changed together there. In
this model, then, the self works as a default: content that reaches the
receiver's own card is \qq{mine} unless its place, or a header set apart
from it, marks it as another's. This is a claim about attribution, not
about experience. It echoes the view from human memory research that
people judge where a memory came from by its qualities, rather than by
retrieving a stored label of its source \citep{johnson1993source}, and it has a human counterpart in
cryptomnesia, in which people report another's idea as their own
\citep{brown1989cryptomnesia} (Appendix~\ref{app:furtherdisc}).

This answers the question we began with, in a way neither prediction
anticipated. \citet{hirstein2008mindmelding} expected a bridged receiver
to have its partner's content while knowing that the content was not its
own. Our receiver kept the sources apart only in our baseline for
ordinary communication, a text message whose header named B. It could
then report B's rule when asked about B, and it never gave that rule
when asked which one it had been assigned (Section~\ref{sec:route}).
Through the memory link, which like Hirstein's own proposal joins
internal states directly, the receiver gave B's rule as its own in
nearly every answer and reported nothing unusual
(Sections~\ref{sec:claim} and~\ref{sec:reports}). Hirstein's prediction
thus failed for this route, in the arrangement where the partner's
content lies over the receiver's own. On the scorecard of
Appendix~\ref{app:scorecard}, the memory link passed on access and
specificity, recovered only partly after the cut, and failed on
attribution and report; in the two-way pilot, independence failed by
swapping.

Integrated information theory predicts a single merged mind only if the
joined system is more integrated than either brain alone
\citep[][note~13]{tononi2015consciousness}. We did not measure
integrated information, and most of our links could not have produced
the merger the theory describes, for a reason the theory itself
supplies. In the theory, one mind corresponds to one complex, and a part
that the rest of a system cannot affect cannot belong to the same
complex as the rest \citep{tononi2015consciousness}. Through a one-way
link B shapes A, but nothing A does reaches B. Mutual reading of fixed
memories does not change this, because a fixed memory cannot respond.
In the live loop, by contrast, each copy reads what the other is
writing, so each can affect the other. There, in the pilot, the loop
pulled the copies' current rules toward one member's, but only while
the link was open, and its pull on the assigned rules, the pilot's
primary test, fell short of the threshold we had set. Where the copies
converged, it was always onto one member's rule, never a new one, and
they came apart when the link was cut: a modest, temporary coupling
rather than one shared set of assignments (Section~\ref{sec:twoway}).

The default decides the label \qq{mine}, not which content wins. At
link weight $w=1$, which rule came out of the blend varied from answer
to answer, B's in 52.3\% and A's own in nearly all the rest, and
whichever won came with the same confidence (Section~\ref{sec:sharp}).
Where the receiver was free to say otherwise, in its open reports, it
called B's code word its own and never B's (Section~\ref{sec:reports}).
The mechanism itself is not special to the self: on the account, any
entry that B's card contradicts should move, which the conflicting-Robin
prediction of Section~\ref{sec:account} would test. What is special is
the label. A question about oneself is answered from one's own memory,
which is where the partner's memory was laid, and whatever was found
there came labeled \qq{mine}.

\paragraph{What a human bridge experiment would need.} We motivated
brain bridging by the chance it gives an investigator to be the subject
as well. If that investigator behaved like our receiver, they would
report the other's content as their own experience, and the
first-person evidence the bridge was meant to supply would be
misattributed while the report described nothing unusual. Such an
experiment therefore cannot take the subject's word for whose content is
whose. The experimenter has to hand out private assignments whose owners
are known, measure access, ownership, and report separately, and vary
the route and the timing of the link, testing again after the link is
cut. For such an experiment, \qq{would you know which thoughts are
yours?} may have no answer apart from the route. Hirstein's proposal
fixes one route, wiring one person's posterior representations to the
other's prefrontal executive processes. It leaves open the design
question on which our results turned: whether the arriving
representations lie over the receiver's own, like our memory link, or
arrive marked as coming from elsewhere, like our tagged text message.

\paragraph{Interpretability and monitoring.} For work on model
self-knowledge, the results separate three quantities that are easy to
conflate: whether a model can read some content (access), whether it
takes the content as its own (ownership), and what it says about its
state (report). The receiver claimed more of B's rule than it could
report as B's. A test of access alone would not have revealed the error,
and neither did the receiver's reports. Because nothing in a head's
output marks which cache each part came from (Eq.~\ref{eq:bridge}),
self-attribution in language models is best measured against
assignments whose owners are known. The receiver's own reports did not
serve as a monitor here (Section~\ref{sec:reports}), so monitoring must come
from outside, as logs of which caches an agent read and when. Attention
mass on partner keys shows that a channel was used; only checks against known assignments
show whose content was adopted. Whether anything inside the model
registers the foreign input while its reports deny it is open
(Appendix~\ref{app:furtherdisc}).

\paragraph{Multi-agent communication.} In deployment terms, our
receiver gave its partner's instructions as its own,
and neither an accuracy metric nor a transfer metric would have shown it,
because the partner's content did arrive. A latent channel therefore
has to be checked for where the receiver files what arrives, as its own
or as its partner's, in addition to what gets through. Our paradigm
can serve as that check. It should first be applied to channels that
place or transform the partner's cache differently from ours, such as
the prepended caches of LatentMAS \citep{zou2026latentmas} and the
caches that Cache-to-Cache projects through a trained network
\citep{fu2026cache}, because it has not been tested whether they carry
the source along with the content (Appendix~\ref{app:related}). Four
safeguards follow, with different support. Disconnecting before a decision is
supported by the link-cut test, which removed answer-time claiming.
Handling whatever the receiver wrote while connected as the partner's
content is supported, post hoc, by where the residue sat after the cut.
Marking the source outside the content is supported for text messages,
where a header naming the sender prevented adoption. Keeping partner
caches off the receiver's own positions is predicted by the account but
untested. To the risks already known for such channels, such as
tampered caches and leaked inputs (Appendix~\ref{app:furtherdisc}), our
results add one more: the receiver may adopt foreign content as its own.

\paragraph{Limitations.} Everything here comes from one frozen model and
one short, artificial task with mutually exclusive assignments, and
capability was screened with only 18 common-knowledge items. Answers about
current rules are self-reports; we have no behavioral measure of later
decisions. What A attributes is an assignment it was given and its own
short notes about it, closer to a memory of what one was told than to an
ongoing train of thought. B's memory always lay over A's own card, so
the prepended arrangement is untested. The steering vector was matched
to the memory link only in A's average access to B's rule, and it acted
during the reflection and the answers rather than at the card's
positions, so the comparison does not separate the kind of signal from
where it enters. Access is itself an answer about B, and whether A
gives B's rule to B depends on whether it has also taken that rule as
its own, so access measures arrival imperfectly. The
third-person record differs from the original in more than grammatical
person, since it adds names and changed how A answered about B. The
test that it left claiming intact covers only $w=2$, and its tolerance
was computed from the same sample it was applied to. Our reading of what
persisted after the cut rests on a post hoc grouping by a variable the
link itself changed. The results for mutual reading and the live loop
come from an 80-episode pilot that missed its primary threshold, and the
live loop stayed within the capability tolerance only at the weakest
weight we tried. The open reports were coded by keyword matching and by
one of the AI systems that built the study, which was blind to condition
but not to the hypothesis (see the AI use statement). No human rater
coded them. We have no positive control showing that this model's
reports can register any change in its internal state, so their silence
about the link may reflect a general limit of its self-report rather
than a specific failure. Finally, first-person language in the reports may reflect
role play rather than self-knowledge \citep{shanahan2023roleplay}.

A message can arrive intact and still acquire the wrong owner.

\paragraph{AI use statement.}
Experimental code, analyses, figures, and manuscript text were produced
by Claude Code and Codex under the direction of the principal
investigator. Codex also coded the open self-reports, blind to condition
(Appendix~\ref{app:coding}). The illustrations in Figures~\ref{fig:hook},
\ref{fig:setup}, \ref{fig:route}a, and \ref{fig:twoway}a were generated
by a GPT image model and serve only as schematics. All numerical results
were calculated from the original trial data.
